% TC-0009 test-case definition (WSP test-case template).
\documentclass[letterpaper]{article}

\usepackage[letterpaper,margin=1in]{geometry}

% Compile from the repository root so this path resolves to the pinned WSP
% submodule.
\input{wsp/testing/test-case-library.tex}

\begin{document}

\def\TCID{TC-0009}
\def\TCTitle{Corpus 213: UpDown buddy control}
\def\TCPurpose{%
  Verify that the buddy of an UpDown (a text box) follows Value, that Value reads the number the buddy holds kept inside Min and Max, that a buddy without a number leaves Value as it is, that AutoBuddy finds the preceding control, and that SyncBuddy False stops the following.}
\def\TCPriority{\TCPriorityHigh}

\def\TCRequirementRef{REQ-0001}
\def\TCDesignRef{Not applicable.}
\def\TCFeatureRef{\texttt{tests/corpus/213-mscomctl2-updown-buddy.vbp}, \texttt{.frm}, \texttt{.ui}, \texttt{.debug} and \texttt{.expected}.}
\def\TCTestTechnique{%
  Scenario test: a program that prints what it observes, compared with a stored transcript.}

\def\TCPreconditions{%
  A \texttt{wfc} executable that includes this library (WFC version 0.8.28 or later) and the files of the case in \texttt{tests/corpus}.}
\def\TCEnvironment{%
  Windows x64. The \texttt{wfc} executable of WFC, given to CMake as \texttt{WFC\_EXE}; the program is run by the interpreter (\texttt{--interpret}) with the system common controls. The case is registered only when \texttt{WFC\_EXE} is set; without it CTest reports the case as skipped.}
\def\TCAssumptions{%
  Not applicable.}

\def\TCInputData{%
  Text1 set to 8, 99 and abc; UpDown1 Value 2; UpDown2 Value 20 and 30 with SyncBuddy off; the script lines \texttt{up} and \texttt{down} on UpDown1 and \texttt{up} on UpDown2.}
\def\TCInitialState{%
  No output of an earlier run is used. CTest test name: \texttt{TC-0009-corpus-updown-buddy}.}

\def\TCProcedure{%
  \begin{enumerate}
    \item Run \texttt{wfc --interpret tests/corpus/213-mscomctl2-updown-buddy.vbp}. Form\_Load prints the first values of both pairs, types 8, 99 and abc into Text1 and prints Value each time, then sets the values of both UpDown controls.
    \item The script clicks the arrows of UpDown1 and the up arrow of UpDown2, then closes the form.
    \item Compare the Immediate-window output with \texttt{213-mscomctl2-updown-buddy.debug} and standard output with the empty \texttt{213-mscomctl2-updown-buddy.expected}.
  \end{enumerate}}
\def\TCExpectedResult{%
  Typing 8 gives Value 8, 99 is kept at 9 (Max), abc leaves 9; the text shows 2 after Value 2; the AutoBuddy text shows 20; after SyncBuddy is False the text stays 20 when Value is 30. Both outputs equal the stored files.}
\def\TCPostConditions{%
  No files are written; the interpreter's windows are closed by the script.}

\TCTemplate

\end{document}
